%! Author = andrei
%! Date = 3/2/26

% Preamble
HADDOCK stands for High Ambiguity Driven protein-protein DOCKing, is a
widely used computational tool for the integrative modelling of biomolecular
interactions. It integrates various types of experimental data, biochemical,
biophysical, bioinformatic prediction, and knowledge  to guide the docking
process.

Haddpcl harnesses the power of CNS (Crystallography and NMR System) for
structure calculation of molecular complexes. What distinguishes Haddock from
other docking software is its ability, inherited from CNS, to incorporate
experimental data as restraints and use these to guide the docking process
alongside traditional energetics and shape complementarity. Moreover, the
intimate coupling with CNS endows Haddock with the ability to actually produce
models of sufficient quality to be archived in the Protein Data Bank.

A central aspect of HADDOCK is the definition of Ambiguous Interaction
Restraints or AIRs. These allow the translation of raw data such as NMR
chemical shift perturbation or mutagenesis experiments into distance
restraints that are incorporated into the energy function used in the calculations.
AIRs are defined through a list of residues that fall under two categories:
active and passive. Generally, active residues are those of central importance
for the interaction, such as residues whose knockouts abolish the interaction or
those where the chemical shift perturbation is higher. Throughout the simulation,
these active residues are restrained to be part of the interface, if possible,
otherwise incurring a scoring penalty. Passive residues are those that contribute
to the interaction but are deemed of less importance. If such a residue does
not belong in the interface there is no scoring penalty. Hence, a careful
selectoin of which residues are active and which are passive is critical for the
success of the docking.

\subsection{HADDOCK scoring function}
CNS modules use the HADDOCK scoring function to score and rank generated models.
The HADDOCK scoring function consists of a linear combination of various weighted
physics-based energy terms and buried surface area.

The scoring is performed according to the weighted sum (HADDOCK score) of the
following 6 terms:
\begin{itemize}
    \item Eelec: electrostatic intermolecular energy
    \item Evdw: van der Waals intermolecular energy
    \item Edesol: desolvation energy
    \item BSA: buried surface area
    \item Eair: distance restraints energy (only unambigous and AIR (ambig) restraints)
    \item Esym: symmetry restraints energy (NCS and C2/C3/C5 terms)
\end{itemize}

As the weights for each of the scoring function components differs for the various
available CNS module, they will be described in each of the modules. These weights
can be tuned by the user, by modifying their related parameters:
\begin{itemize}
    \item \verb|w_elec|: to tune the electrostatic intermolecular energy weight
    \item \verb|w_vdw|: to tune the van der Waals intermolecular energy weight
    \item \verb|w_desolv|: to tune the desolvation energy weight
    \item \verb|w_bsa|: to tune the buried surface area weight
    \item \verb|w_air|: to tune the distance restraints energy weight
    \item \verb|w_sym|: to tune the symmetry restraints energy (NCS and C2/C3/C5 terms) weight
\end{itemize}

